\documentclass{ctxdoc}
\usepackage{ulem}
%\usepackage[markup=underlined,authormarkup=none,commentmarkup=todo]{changes}
\begin{document}
    \title{The \textsf{GZUthesis} package}
    \author{https://github.com/francisol/GZUthesis-template/}
    \maketitle

    \section*{更新历史}

    \begin{itemize}
        \item[1.0.0] (2022-08-17) 初始版本
        \begin{itemize}
            \item 间距格式调整
            \item 内置华文行楷
            \item 输出版本控制(完成部分)
            \item 兼容Mac/Linux
            \item 兼容 overleaf
            \item 答辩委员会名单
            \item 附录
            \item 摘要关键字
            \item 目录
        \end{itemize}

        \item[1.1.0] (2024-08-18) 重构
        \begin{itemize}
            \item 规避版权风险，删除了内置字体
            \item 增加代码注释
            \item 参考文献设置移除宏包
        \end{itemize}
        \item[1.1.1] (2024-08-19) 添加外部字体支持
        \begin{itemize}
            \item 添加外部字体支持
        \end{itemize}
        \item[1.1.2] (2025-01-02) 调整字体搜索策略，以及其他优化
        \begin{itemize}
            \item 增加模板参数，调整 \texttt{Times New Roman},\texttt{ Arial}, \texttt{华文行楷}字体的搜索策略，优化编译
            \item 其他优化
        \end{itemize}
      \item[1.1.3] (2026-03-13) 添加双面打印空白页水印，优化行间距
      \item[1.1.4] (2026-04-22) 调整论文模板排版与校徽资源，添加盲审回复模版
      \item[2.0.0] (2026-04-22) 根据贵州大学研究生院 2026 年 3 月 9 日发布的
      \href{https://gs.gzu.edu.cn/2026/0309/c11824a265403/page.htm}{《关于规范贵州大学研究生学位论文写作格式的通知》}
      调整新版论文格式与字体策略
      \item[2.1.0] (2026-07-22) 承诺书自动排版与答辩委员会名单单独打印
      \begin{itemize}
          \item 新增 \texttt{metadata.tex} 集中管理论文元数据，主论文与答辩委员会名单入口共用
          \item 终稿封面后自动生成承诺书，支持签名图片叠加、日期与保密/不保密勾选
          \item 新增 \texttt{defense\_committee.tex} 入口与 \cs{makedefensepage}，可单独打印答辩委员会名单
          \item 新增 \cs{defensepdf}，支持直接插入已签字的答辩委员会名单扫描件
      \end{itemize}
        
    \end{itemize}


    \section{说明}
    \textbf{该文档为用户手册，并非宏包使用实例。}


    本模板提供贵州大学硕士博士\sout{以及本科生版本的}学位(毕业)论文模板。
    本模板当前版主要依据贵州大学研究生院 2026 年 3 月 9 日发布的
    \href{https://gs.gzu.edu.cn/2026/0309/c11824a265403/page.htm}{《关于规范贵州大学研究生学位论文写作格式的通知》}
    及其附件《贵州大学研究生学位论文格式》调整；旧版模板依据 2022 年发布《贵州大学研究生学位论文格式》、《贵州大学研究生学位论文盲评封面格式》制作。

    本模板亦在\href{https://www.overleaf.com/latex/templates/gui-zhou-da-xue-yan-jiu-sheng-bi-ye-lun-wen-mo-ban/tstnysppbxvh}{Overleaf} 提供

    如有疑问请在 \href{https://github.com/francisol/GZUthesis-template/issues}{Github} 上提交 issues。

    \section{内置信息}
    \subsection{内置宏包}
    \begin{description}[align=left,leftmargin=!,labelwidth=5em]
        \item[xkeyval] 解析宏包参数，用于辅助编写宏包
        \item[ifthen] 用于辅助编写宏包
        \item[ctex] 中文字符处理
        \item[geometry] 设置页边距
        \item[hyperref] 链接
        \item[appendix] 实现附录
        \item[setspace] 设置行间距
        \item[titletoc] 设置目录
        \item[fancyhdr] 页面布局，实现页眉页脚
        \item[pdfpages] 插入 PDF 文件，实现插入已签字答辩委员会名单
        \item[ulem] 用于实现下划线样式
        \item[bicaption] 用于实现中英图表说明
        \item[draftwatermark] 开启空白页水印时启用
        \item[xcolor] 开启空白页水印时启用
        \item[graphicx] 插入图片、承诺书底图和签名图片
        \item[tikz] 在承诺书底图上叠加签名、日期和保密勾选
    \end{description}

    \subsection{内置样式}
    已增加中英文图表说明，图表默认五号黑体。

    参考文献样式未在宏包中限制，但在样例文件中给出国标格式。

    \section{已知问题}
    默认中文字体由 \textbf{ctex} 选择。Windows 系统通常不需要开启 \textbf{windows} 选项；该选项主要用于在非 Windows 系统上显式使用已安装的 Windows 常用中文字体，或希望固定使用 \texttt{SimSun}/\texttt{FangSong}/\texttt{SimHei}/\texttt{KaiTi} 等字体时。
    若只想单独启用 \texttt{Times New Roman}, \texttt{Arial} 字体，请使用包选项 \textbf{tnra}。

    启用 \textbf{windows} 选项后，模板会先查找系统字体名和当前目录字体文件；若仍未找到，会读取 \textbf{ctex}/\textbf{xeCJK} 已定义的默认字体族作为兜底。因此即使系统缺少某个 Windows 字体，也会保留可用的系统默认字体。

    盲审版首页 \textbf{XX 届X士研究生学位论文} 字样为华文行楷（STXINGKA）。使用包选项 \textbf{windows} 时会自动尝试启用华文行楷；若只想单独启用该字体，可使用包选项 \textbf{hwxk}。


    \textbf{页面布局并未一比一复刻，但是没那么离谱。}


    \section{使用}
    \subsection{推荐文件组织}

    模板仓库提供了多个入口文件。一般只需要修改 \texttt{metadata.tex} 和各章节文件，不建议把论文元数据散落在多个入口中。

    \begin{description}
        \item[\texttt{metadata.tex}] 论文元数据文件，集中填写题名、作者、专业、导师、答辩委员会、承诺书日期和签名图片等信息。主论文和单独答辩委员会名单入口都会读取它。
        \item[\texttt{final\_thesis.tex}] 主论文入口。默认读取 \texttt{metadata.tex}，再生成封面、前置部分、正文、参考文献和附录。
        \item[\texttt{defense\_committee.tex}] 答辩委员会名单单独打印入口。它读取同一份 \texttt{metadata.tex}，只输出答辩委员会名单页。
        \item[\texttt{chapter/}] 章节文件目录。摘要、绪论、致谢和附录示例均放在这里。
        \item[\texttt{reference.bib}] 参考文献数据库。
        \item[\texttt{logo/}] 校徽、承诺书底图等模板资源目录。签名图片也可以放在这里。
        \item[\texttt{out/}] 编译输出目录，由 \texttt{latexmk} 自动生成。
    \end{description}

    推荐的修改顺序如下：
    \begin{enumerate}
        \item 修改 \texttt{metadata.tex} 中的论文题名、作者、专业、导师和答辩委员会信息。
        \item 在 \texttt{chapter/abstract.tex}、\texttt{chapter/intro.tex} 等章节文件中写正文。
        \item 在 \texttt{reference.bib} 中维护参考文献。
        \item 用 \texttt{latexmk} 编译主论文；若只需要打印答辩委员会名单，则用 \texttt{latexmk defense\_committee.tex}。
    \end{enumerate}

    \subsection{编译方式}

    本模板使用 XeLaTeX 编译，推荐直接运行：
    \begin{verbatim}
latexmk
    \end{verbatim}
    该命令会读取 \texttt{.latexmkrc}，默认编译 \texttt{final\_thesis.tex}，并把中间文件和 PDF 输出到 \texttt{out/} 目录。

    常用命令如下：
    \begin{description}
        \item[\texttt{latexmk}] 编译默认主论文入口，输出 \texttt{out/final\_thesis.pdf}。
        \item[\texttt{latexmk final\_thesis.tex}] 显式编译主论文。
        \item[\texttt{latexmk defense\_committee.tex}] 单独编译答辩委员会名单，输出 \texttt{out/defense\_committee.pdf}。
        \item[\texttt{latexmk response.tex}] 编译盲审意见回复示例。
        \item[\texttt{latexmk -c}] 清理辅助文件，保留已经生成的 PDF。
    \end{description}

    \subsection{最小主论文入口}

    一个典型的主论文入口如下。实际模板中的 \texttt{final\_thesis.tex} 已经采用这种结构。

    \begin{verbatim}
\documentclass[oneside,degree=master,print=final,windows]{GZUthesis}
\usepackage{gbt7714}
\usepackage{booktabs}

% 论文元数据示例。正式写作时优先修改本文件，主论文、答辩委员会名单等入口会复用这里的信息。
\gradyear{2021}
\major{学科专业}
\research{研究方向}
\author{学生}
\supervisor{导师}
\stuid{(学号)}
\title{深度神经网络在图像像识别中的应用与优化识别中的应用与优化的应用与优化识别中的应用与优化}
\entitle{How are You How are You How are You How are You How are You How are You How are You How are You How are You}
\department{培养单位}
\seclvl{<密级>}
\classifyno{<分类号>}
\udcno{<UDC分类号>}
\degreecategory{学术/专业学位}
\date{2026年7月22日}

% 答辩委员会名单信息。
\chairman{贵州大学}{张XX}{教授}
\addexpert{贵州没有大学}{李XX}{教授}
\addexpert{贵州大学}{李XX}{教授}
\addexpert{贵州大学}{李XX}{超级教授}
\addexpert{贵州大学}{刘XX}{服教授}
\defsec{张三\ 学士}
\defaddr{兜率宫}

% 如需在终稿中使用签字后的答辩委员会名单扫描件，取消下一行注释。
% \defensepdf{figures/defense-committee-signed.pdf}

% 终稿承诺书签名与日期。图片为空时只输出承诺书底图、日期和保密勾选。
% \studentsignature[2.55cm]{logo/student-signature.png}
% \supervisorsignature[2.55cm]{logo/supervisor-signature.png}
\promisedate{2026}{7}{22}
\promisenonconfidential


\begin{document}
  \maketitle
  \frontmatter
  \input{chapter/abstract.tex}
  \tableofcontents
  \listoffigures
  \listoftables
  \printnomenclature

  \mainmatter
  \input{chapter/intro.tex}

  \bibliographystyle{gbt7714-numerical}
  \bibliography{reference}
\end{document}
    \end{verbatim}

    \subsection{元数据文件示例}

    \texttt{metadata.tex} 应只包含设置命令，不包含 \cs{documentclass} 和 \cs{begin}\{document\}。示例：

    \begin{verbatim}
\gradyear{2026}
\major{软件工程}
\research{软件工程理论与计算理论}
\author{张三}
\supervisor{李四\ 教授}
\stuid{2026010001}
\title{论文中文题目}
\entitle{English Thesis Title}
\department{计算机科学与技术学院}
\seclvl{公开}
\classifyno{TP301}
\udcno{004}
\degreecategory{学术学位}
\date{2026年7月22日}

\chairman{贵州大学}{王五}{教授}
\addexpert{贵州大学}{赵六}{教授}
\addexpert{某某大学}{钱七}{教授}
\defsec{周八\ 副教授/博士}
\defaddr{贵州大学某楼某会议室}

\promisedate{2026}{7}{22}
\promisenonconfidential
    \end{verbatim}

    如果需要把签名图片叠加到承诺书上，可以继续加入：
    \begin{verbatim}
\studentsignature[2.55cm]{logo/student-signature.png}
\supervisorsignature[2.55cm]{logo/supervisor-signature.png}
    \end{verbatim}

    \subsection{输出模式与自动页面}

    \pkg{GZUthesis} 通过 \texttt{print} 选项控制 \cs{maketitle} 生成的页面。

    \begin{description}
        \item[\texttt{print=review}] 生成盲审封面。该模式不输出个人信息、答辩委员会名单和承诺书。
        \item[\texttt{print=final}] 生成终稿封面，并自动追加答辩委员会名单和承诺书。承诺书会使用 \texttt{logo/promise-normalized.pdf} 作为底图，并叠加日期、保密勾选和可选签名图片。
        \item[\texttt{print=draft}] 生成简洁草稿封面，便于写作阶段快速校对正文。
    \end{description}

    终稿模式下不需要在正文末尾手动调用 \cs{makepromisepage}；\cs{maketitle} 已经会自动处理。若手动再调用一次，会重复生成承诺书。

    \subsection{单独打印答辩委员会名单}

    模板提供 \cs{makedefensepage} 用于单独生成答辩委员会名单。仓库中的 \texttt{defense\_committee.tex} 即为示例：

    \begin{verbatim}
\documentclass[oneside,degree=master,print=draft,windows]{GZUthesis}
% 论文元数据示例。正式写作时优先修改本文件，主论文、答辩委员会名单等入口会复用这里的信息。
\gradyear{2021}
\major{学科专业}
\research{研究方向}
\author{学生}
\supervisor{导师}
\stuid{(学号)}
\title{深度神经网络在图像像识别中的应用与优化识别中的应用与优化的应用与优化识别中的应用与优化}
\entitle{How are You How are You How are You How are You How are You How are You How are You How are You How are You}
\department{培养单位}
\seclvl{<密级>}
\classifyno{<分类号>}
\udcno{<UDC分类号>}
\degreecategory{学术/专业学位}
\date{2026年7月22日}

% 答辩委员会名单信息。
\chairman{贵州大学}{张XX}{教授}
\addexpert{贵州没有大学}{李XX}{教授}
\addexpert{贵州大学}{李XX}{教授}
\addexpert{贵州大学}{李XX}{超级教授}
\addexpert{贵州大学}{刘XX}{服教授}
\defsec{张三\ 学士}
\defaddr{兜率宫}

% 如需在终稿中使用签字后的答辩委员会名单扫描件，取消下一行注释。
% \defensepdf{figures/defense-committee-signed.pdf}

% 终稿承诺书签名与日期。图片为空时只输出承诺书底图、日期和保密勾选。
% \studentsignature[2.55cm]{logo/student-signature.png}
% \supervisorsignature[2.55cm]{logo/supervisor-signature.png}
\promisedate{2026}{7}{22}
\promisenonconfidential

\begin{document}
\makedefensepage
\end{document}
    \end{verbatim}

    编译命令：
    \begin{verbatim}
latexmk defense_committee.tex
    \end{verbatim}

    如果已经有签字后的答辩委员会名单扫描件，可以在 \texttt{metadata.tex} 中设置：
    \begin{verbatim}
\defensepdf{figures/defense-committee-signed.pdf}
    \end{verbatim}
    此时 \texttt{print=final} 的主论文和 \cs{makedefensepage} 都会直接插入该 PDF 的第一页，而不是重新排版答辩委员会名单。

    \subsection{承诺书与签名}

    终稿模式会自动插入承诺书。默认情况下，即使没有设置签名图片，也会生成承诺书底图、日期和保密/不保密勾选。常用设置如下：

    \begin{verbatim}
\promisedate{2026}{7}{22}
\promisenonconfidential

% 或者：保密 3 年
\promiseconfidential{3}

% 可选：叠加签名图片
\studentsignature[2.55cm]{logo/student-signature.png}
\supervisorsignature[2.55cm]{logo/supervisor-signature.png}
    \end{verbatim}

    如果签名位置需要微调，可以使用：
    \begin{verbatim}
\promiseoriginalitysignatureposition{9.0cm}{18.65cm}
\promiseauthorizationstudentposition{8.90cm}{5.48cm}
\promiseauthorizationsupervisorposition{14.25cm}{4.83cm}
    \end{verbatim}

    也可以使用集中配置接口：
    \begin{verbatim}
\promisesetup{
  student-signature=logo/student-signature.png,
  supervisor-signature=logo/supervisor-signature.png,
  student-width=2.55cm,
  supervisor-width=2.55cm,
  date-year=2026,
  date-month=7,
  date-day=22,
  confidential=false
}
    \end{verbatim}

    \subsection{宏包选项}
    \begin{function}[added=2024-08-18]{degree}
        \begin{syntax}
            degree = <doctor|master|bachelor>
        \end{syntax}
        指定该文档的学位类型，默认为博士论文（doctor）。该值仅影响文档的标题。
    \end{function}
    \begin{function}[added=2024-08-18]{print}
        \begin{syntax}
            print = <review|final|draft>
        \end{syntax}
        指定该文档的打印类型，默认为草稿（draft）。该值影响文档的封面以及正文之前的页面展示。
    \end{function}

    \begin{optdesc}
        \item[review] 盲审版本，封面不显示个人信息，封面之后无附加信息。
        \item[final] 最终版本，封面显示个人信息，封面之后自动显示答辩委员会名单和承诺书。
        \item[draft] 草稿版本，用于自己校验，只显示正文。
    \end{optdesc}
    
    
    \begin{function}[added=2024-08-19]{windows}
        \begin{syntax}
            windows
        \end{syntax}
        若使用 windows 参数，则尝试在系统或者文档当前目录下搜索 Windows 常用论文排版字体并使用，同时自动启用 \textbf{hwxk} 与 \textbf{tnra} 字体检查。Windows 系统通常不需要开启该选项，因为 \textbf{ctex} 会自动选择系统中文字体；该选项更适合非 Windows 系统显式启用已安装的 Windows 字体。
        \begin{description}
            \item[宋体] 搜索系统字体 \texttt{SimSun} 或者当前目录下 \texttt{simsun.ttc}；找不到时读取默认 \texttt{zhsong}
            \item[仿宋] 搜索系统字体 \texttt{FangSong} 或者当前目录下 \texttt{simfang.ttf}；找不到时读取默认 \texttt{zhfs}
            \item[黑体] 搜索系统字体 \texttt{SimHei} 或者当前目录下 \texttt{simhei.ttf}；找不到时读取默认 \texttt{zhhei}
            \item[楷体] 搜索系统字体 \texttt{KaiTi} 或者当前目录下 \texttt{simkai.ttf}；找不到时读取默认 \texttt{zhkai}
         \end{description}
         宋体族设置 \texttt{AutoFakeBold=2.5}，并将 \texttt{ItalicFont} 指向找到或兜底得到的楷体，因此中文 \cs{itshape} 会显示为楷体。
    \end{function}

    \begin{function}[added=2025-01-02]{hwxk}
        \begin{syntax}
            hwxk
        \end{syntax}
        若使用 hwxk 参数，则尝试在系统或者文档当前目录下搜索 \texttt{华文行楷} 字体并使用。使用 \textbf{windows} 参数时会自动启用该检查。
        \begin{description}
            \item[华文行楷] 搜索系统字体 \texttt{STXINGKA} 或者当前目录下 \texttt{STXINGKA.TTF}

         \end{description}
    \end{function}

    \begin{function}[added=2025-01-02]{tnra}
        \begin{syntax}
            tnra
        \end{syntax}
        若使用 tnra 参数，则尝试在系统搜索 \texttt{Arial} 和 \texttt{Times New Roman} 字体并使用。使用 \textbf{windows} 参数时会自动启用该检查。
        \begin{description}
            \item[Arial] 搜索系统字体 \texttt{Arial}
            \item[Times New Roman] 搜索系统字体 \texttt{Times New Roman}
         \end{description}
    \end{function}

    \begin{function}[added=2026-03-13]{watermark}
        \begin{syntax}
            watermark
        \end{syntax}
        若使用 watermark 参数，则当双面打印时会在空白页面添加``此页无正文''水印
    \end{function}

    \subsection{宏包命令}

    本节介绍了本包自定义和重定义的命令。

    \subsubsection{信息设置类命令}
    该类命令均保存为\cs{GZU@*}变量

    \begin{function}[added=2024-08-18]{\title}
        \begin{syntax}
            \cs{title}\marg{论文题目} \\
        \end{syntax}
        该命令设置论文题目，题目展示最多三行，题目会写入到pdf文件属性中
        \begin{texnote}
            该命令为重定义命令
        \end{texnote}
    \end{function}

    \begin{function}[added=2024-08-18]{\author}
        \begin{syntax}
            \cs{author}\marg{作者姓名} \\
        \end{syntax}
        该命令设置论文作者，会填充到“研究生”、“答辩人”等表格内容中，作者会写入到pdf文件属性中
        \begin{texnote}
            该命令为重定义命令，不会覆盖\cs{\@author}。
        \end{texnote}

    \end{function}

    \begin{function}[added=2026-04-22]{\department}
        \begin{syntax}
            \cs{department}\marg{培养单位} \\
        \end{syntax}
        该命令设置培养单位，会填充到终稿封面的“培养单位”字段中，默认留空。
    \end{function}

    \begin{function}[added=2024-08-18]{\date}
        \begin{syntax}
            \cs{date}\marg{XXXX年XX月} \\
        \end{syntax}
        该命令设置论文日期，默认为当前时间。日期中的年月用于封面底部，完整日期也会同步作为答辩日期。
        日期格式必须为 “XXXX年XX月” 或“XXXX年XX月xx日”，否则可能会报错或显示错误。
        \begin{texnote}
            该命令为重定义命令，不会覆盖\cs{\@date}。若答辩日期不同于封面日期，可以在 \cs{date} 之后再用 \cs{defdate} 单独覆盖。
        \end{texnote}

    \end{function}


    \begin{function}[added=2024-08-18]{\paperno}
        \begin{syntax}
            \cs{paperno}\marg{论文号} \\
        \end{syntax}
        该命令设置论文编号，默认留空。
    \end{function}


    \begin{function}[added=2024-08-18]{\defsec}
        \begin{syntax}
            \cs{defsec}\marg{XXX+职称/学位} \\
        \end{syntax}
        该命令设置答辩秘书，格式为“XXX+职称/学位”，默认留空。
    \end{function}


    \begin{function}[added=2024-08-18]{\addexpert}
        \begin{syntax}
            \cs{addexpert}\marg{单位}\marg{姓名}\marg{职称} \\
        \end{syntax}
        该命令为答辩专家列表添加一个专家信息，可以重复调用，只会追加信息。
        \begin{texnote}
            示例：
            \begin{verbatim}
\addexpert{贵州大学}{李XX}{教授}
            \end{verbatim}
        \end{texnote}
    \end{function}


    \begin{function}[added=2024-08-18]{\chairman}
        \begin{syntax}
            \cs{chairman}\marg{单位}\marg{姓名}\marg{职称} \\
        \end{syntax}
        该命令设置答辩主席，默认留空。答辩委员会名单中会在该行末尾自动标注“（主席）”。
    \end{function}


    \begin{function}[added=2024-08-18]{\major}
        \begin{syntax}
            \cs{major}\marg{学科专业} \\
        \end{syntax}
        该命令设置学科专业，默认留空。
    \end{function}


    \begin{function}[added=2024-08-18]{\class}
        \begin{syntax}
            \cs{class}\marg{班级} \\
        \end{syntax}
        该命令设置班级，默认留空。
        \begin{texnote}
            本科模板预留参数，研究生无效。
        \end{texnote}
    \end{function}


    \begin{function}[added=2024-08-18]{\college}
        \begin{syntax}
            \cs{college}\marg{学院} \\
        \end{syntax}
        该命令设置学院，默认留空。

        \begin{texnote}
            本科模板预留参数，研究生无效。
        \end{texnote}
    \end{function}


    \begin{function}[added=2024-08-18]{\research}
        \begin{syntax}
            \cs{research}\marg{研究方向} \\
        \end{syntax}
        该命令设置研究方向，默认留空。
    \end{function}


    \begin{function}[added=2024-08-18]{\supervisor}
        \begin{syntax}
            \cs{supervisor}\marg{导师} \\
        \end{syntax}
        该命令设置导师，默认留空。
    \end{function}


    \begin{function}[added=2024-08-18]{\stuid}
        \begin{syntax}
            \cs{stuid}\marg{学号} \
        \end{syntax}
        该命令设置学号，默认留空。

    \end{function}



    \begin{function}[added=2024-08-18]{\defaddr}
        \begin{syntax}
            \cs{defaddr}\marg{答辩地点} \\
        \end{syntax}
        该命令设置答辩地点，默认留空。
    \end{function}

    \begin{function}[added=2026-07-22]{\defensepdf}
        \begin{syntax}
            \cs{defensepdf}\marg{PDF 文件路径}
        \end{syntax}
        使用已签字的答辩委员会名单扫描件。设置后，终稿封面之后和 \cs{makedefensepage} 都会直接插入该 PDF 的第一页。
        \begin{texnote}
            如果想恢复模板自动排版的答辩委员会名单，可使用 \cs{cleardefensepdf}。
        \end{texnote}
    \end{function}

    \begin{function}[added=2026-07-22]{\cleardefensepdf}
        \begin{syntax}
            \cs{cleardefensepdf}
        \end{syntax}
        清除通过 \cs{defensepdf} 设置的扫描件路径，恢复自动排版答辩委员会名单。
    \end{function}


    \begin{function}[added=2024-08-18]{\gradyear}
        \begin{syntax}
            \cs{gradyear}\marg{毕业年份} \
        \end{syntax}
        该命令设置毕业年份，默认为当前年份。
    \end{function}


    \begin{function}[added=2024-08-18]{\defdate}
        \begin{syntax}
            \cs{defdate}\marg{答辩时间} \\
        \end{syntax}
        该命令设置答辩时间，默认与 \cs{date} 相同。若答辩日期和封面日期不同，可单独设置该命令。
    \end{function}





    \begin{function}[added=2024-08-18]{\classifyno}
        \begin{syntax}
            \cs{classifyno}\marg{分类号}
        \end{syntax}
        该命令设置分中图类号，默认留空。
    \end{function}

    \begin{function}[added=2026-04-22]{\udcno}
        \begin{syntax}
            \cs{udcno}\marg{UDC分类号}
        \end{syntax}
        该命令设置 UDC 分类号，会填充到终稿封面顶部的“UDC 分类号”字段中，默认留空。
    \end{function}

    \begin{function}[added=2026-04-22]{\degreecategory}
        \begin{syntax}
            \cs{degreecategory}\marg{学位类别}
        \end{syntax}
        设置终稿封面中论文类型下方的学位类别，例如“学术学位”或“专业学位”。
    \end{function}

    \begin{function}[added=2026-04-22]{\entitle}
        \begin{syntax}
            \cs{entitle}\marg{英文题目}
        \end{syntax}
        设置英文题目。终稿封面会在中文题目下方显示英文题目，英文摘要标题也会使用该值。
    \end{function}

    \begin{function}[added=2026-07-22]{\studentsignature}
        \begin{syntax}
            \cs{studentsignature}\oarg{宽度}\marg{图片路径}
        \end{syntax}
        设置承诺书中的学生签名图片。可选宽度默认是 \texttt{2.55cm}。
    \end{function}

    \begin{function}[added=2026-07-22]{\supervisorsignature}
        \begin{syntax}
            \cs{supervisorsignature}\oarg{宽度}\marg{图片路径}
        \end{syntax}
        设置承诺书中的导师签名图片。可选宽度默认是 \texttt{2.55cm}。\cs{advisorsignature} 是该命令的同义命令。
    \end{function}

    \begin{function}[added=2026-07-22]{\promisedate}
        \begin{syntax}
            \cs{promisedate}\marg{年}\marg{月}\marg{日}
        \end{syntax}
        设置承诺书中原创性声明和授权声明共用的签署日期。
    \end{function}

    \begin{function}[added=2026-07-22]{\promisenonconfidential}
        \begin{syntax}
            \cs{promisenonconfidential}
        \end{syntax}
        在承诺书中勾选“不保密”。这是模板示例的默认设置。
    \end{function}

    \begin{function}[added=2026-07-22]{\promiseconfidential}
        \begin{syntax}
            \cs{promiseconfidential}\marg{保密年限}
        \end{syntax}
        在承诺书中勾选“保密”，并填写“在若干年解密后适用授权”中的年限。
    \end{function}

    \begin{function}[added=2026-07-22]{\promisesetup}
        \begin{syntax}
            \cs{promisesetup}\marg{key=value 列表}
        \end{syntax}
        集中配置承诺书签名、日期、保密状态和坐标。常用 key 如下：
        \begin{description}
            \item[\texttt{student-signature}] 学生签名图片路径。
            \item[\texttt{supervisor-signature}] 导师签名图片路径，\texttt{advisor-signature} 为同义 key。
            \item[\texttt{student-width}] 学生签名图片宽度。
            \item[\texttt{supervisor-width}] 导师签名图片宽度，\texttt{advisor-width} 为同义 key。
            \item[\texttt{date-year/date-month/date-day}] 承诺书签署日期。
            \item[\texttt{confidential}] 是否保密，取值为 \texttt{true} 或 \texttt{false}。
            \item[\texttt{secret-years}] 保密年限，仅保密时使用。
            \item[\texttt{originality-student-x/originality-student-y}] 原创性声明页学生签名中心位置（相对页面左下角，cm）。
            \item[\texttt{authorization-student-x/authorization-student-y}] 授权声明页学生签名中心位置（相对页面左下角，cm）。
            \item[\texttt{authorization-supervisor-x/authorization-supervisor-y}] 授权声明页导师签名中心位置（相对页面左下角，cm）。
            \item[\texttt{originality-date-y}] 原创性声明页签署日期基线高度（cm）。
            \item[\texttt{authorization-date-y}] 授权声明页签署日期基线高度（cm）。
        \end{description}
    \end{function}

    \begin{function}[added=2026-07-22]{\promiseoriginalitysignatureposition}
        \begin{syntax}
            \cs{promiseoriginalitysignatureposition}\marg{x}\marg{y}
        \end{syntax}
        调整原创性声明页学生签名图片中心位置，坐标为相对页面左下角的厘米数。默认 \texttt{(9.0, 18.65)}。
    \end{function}

    \begin{function}[added=2026-07-22]{\promiseauthorizationstudentposition}
        \begin{syntax}
            \cs{promiseauthorizationstudentposition}\marg{x}\marg{y}
        \end{syntax}
        调整授权声明页学生签名图片中心位置，坐标为相对页面左下角的厘米数。默认 \texttt{(8.90, 5.48)}。
    \end{function}

    \begin{function}[added=2026-07-22]{\promiseauthorizationsupervisorposition}
        \begin{syntax}
            \cs{promiseauthorizationsupervisorposition}\marg{x}\marg{y}
        \end{syntax}
        调整授权声明页导师签名图片中心位置，坐标为相对页面左下角的厘米数。默认 \texttt{(14.25, 4.83)}。
    \end{function}

    \subsubsection{控制类命令}

    \begin{function}[added=2024-08-18]{\maketitle}
        \begin{syntax}
            \cs{maketitle}
        \end{syntax}
        根据本文前述信息生成封面和前置附加页。具体输出由 \texttt{print} 选项决定：盲审模式只生成盲审封面，终稿模式生成终稿封面、答辩委员会名单和承诺书，草稿模式生成简洁草稿封面。
    \end{function}

    \begin{function}[added=2026-07-22]{\makedefensepage}
        \begin{syntax}
            \cs{makedefensepage}
        \end{syntax}
        单独生成答辩委员会名单。推荐在 \texttt{defense\_committee.tex} 中使用，并与主论文共用 \texttt{metadata.tex}。
    \end{function}

    \begin{function}[added=2026-07-22]{\makepromisepage}
        \begin{syntax}
            \cs{makepromisepage}
        \end{syntax}
        手动生成承诺书页。一般不需要直接调用，因为 \texttt{print=final} 下 \cs{maketitle} 会自动插入承诺书。
    \end{function}

    \begin{function}[added=2024-08-18]{\frontmatter}
        \begin{syntax}
            \cs{frontmatter}
        \end{syntax}
        正文前样式，无页眉，页脚只保留罗马数字的页码信息
    \end{function}

    \begin{function}[added=2024-08-18]{\mainmatter}
        \begin{syntax}
            \cs{mainmatter}
        \end{syntax}
        正文样式，页眉含有章节信息以及论文题目，页脚只保留阿拉伯数字的页码信息。
    \end{function}


    \begin{function}[added=2024-08-18]{\keywords}
        \begin{syntax}
            \cs{keywords}\{人工智能, 机器学习\}\\
            \cs{keywords}[Key words:]\{AI, ML\}
        \end{syntax}
        按要求格式输出关键字信息，接收两个参数，第一个参数为关键词标题,默认为 \textbf{关键字：};第二个为关键词具体内容。
    \end{function}



    \begin{function}[added=2024-08-18]{\acknowledgments}
        \begin{syntax}
            \cs{acknowledgments}\\
            \cs{acknowledgments}\marg{致谢}
        \end{syntax}
        生成致谢章节，章节不编码，标题居中。命令参数为章节标题，默认为 \textbf{致谢}
    \end{function}

    \subsubsection{缩略语、术语与符号表命令}

    模板内置了基于 \pkg{nomencl} 的“缩略语、术语、符号对照表”。在正文开始前集中登记条目，并在目录之后、正文之前使用 \cs{printnomenclature} 输出。

    \begin{function}[added=2026-04-29]{\abbrdef}
        \begin{syntax}
            \cs{abbrdef}\marg{中文名称}\marg{英文全称}\marg{缩写}
        \end{syntax}
        登记缩略语条目，固定写入 A 组“缩略语对照表”，不在当前位置输出正文。
        \begin{texnote}
            示例：
            \begin{verbatim}
\abbrdef{卷积神经网络}{Convolutional Neural Network}{CNN}
\abbrdef{生成对抗网络}{Generative Adversarial Network}{GAN}
            \end{verbatim}
        \end{texnote}
    \end{function}

    \begin{function}[added=2026-04-29]{\abbritem}
        \begin{syntax}
            \cs{abbritem}\marg{中文名称}\marg{英文全称}\marg{缩写}
        \end{syntax}
        在当前位置输出“中文名称（英文全称，缩写）”，同时登记到缩略语表。
    \end{function}

    \begin{function}[added=2026-04-29]{\nomdef}
        \begin{syntax}
            \cs{nomdef}\marg{中文术语}\marg{英文术语}
        \end{syntax}
        登记双语术语条目，固定写入 N 组“术语表”，不在当前位置输出正文。
        \begin{texnote}
            示例：
            \begin{verbatim}
\nomdef{深度学习}{Deep Learning}
            \end{verbatim}
        \end{texnote}
    \end{function}

    \begin{function}[added=2026-04-29]{\nomitem}
        \begin{syntax}
            \cs{nomitem}\marg{中文术语}\marg{英文术语}
        \end{syntax}
        在当前位置输出“中文术语（英文术语）”，同时登记到术语表。
    \end{function}

    \begin{function}[added=2026-04-29]{\termitem}
        \begin{syntax}
            \cs{termitem}\marg{中文术语}\marg{解释}
        \end{syntax}
        在当前位置输出“中文术语”，同时登记到术语表。
    \end{function}

    \begin{function}[added=2026-04-29]{\symbolitem}
        \begin{syntax}
            \cs{symbolitem}\marg{符号}\marg{说明}
        \end{syntax}
        登记符号条目，固定写入 S 组“符号表”，不在当前位置输出正文。
        \begin{texnote}
            示例：
            \begin{verbatim}
\symbolitem{$\alpha$}{学习率，用于控制模型参数每次更新的步长。}
\symbolitem{$\mathcal{D}$}{训练数据集。}
            \end{verbatim}
        \end{texnote}
    \end{function}

    \begin{function}[added=2026-04-29]{\printnomenclature}
        \begin{syntax}
            \cs{printnomenclature}
        \end{syntax}
        输出“缩略语、术语、符号对照表”。推荐放在 \cs{tableofcontents} 之后、\cs{mainmatter} 之前。
        \begin{texnote}
            使用 \texttt{latexmk} 编译时会自动调用 \texttt{makeindex} 生成对照表；手动编译时需要在 XeLaTeX 之后运行：
            \begin{verbatim}
makeindex out/final_thesis.nlo -s nomencl.ist -o out/final_thesis.nls
            \end{verbatim}
        \end{texnote}
    \end{function}

    \subsection{宏包环境}

    \begin{environment}{abstract}
        \begin{syntax}
            \cs{begin}\{abstract\}[abstract] \\
                This research investigates the latest applications and optimization techniques
            of deep neural networks in the field of image recognition. \\
            \cs{end}\{abstract\}
        \end{syntax}
        这个环境用于创建摘要章节，章节名称通过可选参数指定，默认为 \textbf{摘要}。
        \begin{texnote}
            abstract环境在book中不存在。
        \end{texnote}
    \end{environment}


    \begin{environment}{appendices}
        \begin{syntax}
            \cs{begin}\{appendices\} \\
            \quad  \cs{chapter}\{参与会议\} \\
            \cs{end}\{appendices\}
        \end{syntax}
        这个环境用于创建附录章节，章节不编码，包裹在其中的章节按大写英文字母编码。

        \begin{texnote}
            使用示例：
            \begin{verbatim}
            \begin{appendices}
                \chapter{参与会议}
            \end{appendices}
            \end{verbatim}

            目录会新增两个一级目录分别为 \textbf{附录} 、\textbf{附录 A 参与会议}。

        \end{texnote}

    \end{environment}

\end{document}
